\documentclass[11pt]{article}
\usepackage[margin=1in]{geometry}
\usepackage{amsmath,amssymb,amsthm,mathtools}
\usepackage[T1]{fontenc}
\usepackage{lmodern,microtype}
\usepackage[hidelinks]{hyperref}
\hypersetup{pdftitle={Polynomial-Time Exact Learning of Fourier-Sparse Boolean Functions from Nearly Linear Quantum Samples},pdfauthor={Daniel Fredriksen}}
\newtheorem{theorem}{Theorem}
\newtheorem{lemma}[theorem]{Lemma}
\newtheorem{proposition}[theorem]{Proposition}
\theoremstyle{remark}\newtheorem{remark}[theorem]{Remark}
\newcommand{\F}{\mathbb F}
\newcommand{\E}{\mathbb E}
\newcommand{\supp}{\operatorname{supp}}
\newcommand{\sgn}{\operatorname{sign}}
\newcommand{\FAIL}{\mathsf{FAIL}}
\newcommand{\norm}[1]{\left\|#1\right\|}
\title{Polynomial-Time Exact Learning of Fourier-Sparse Boolean Functions from Nearly Linear Quantum Samples}
\author{Daniel Fredriksen\\\small Quantyra Inc.\\\small\href{https://orcid.org/0009-0002-9123-7435}{ORCID: 0009-0002-9123-7435}}
\date{November 10, 2026\\\small Manuscript version 0.4.0}
\begin{document}
\maketitle
\begin{abstract}
We give an algorithm that exactly recovers the Fourier representation of every Boolean function with at most $k$ nonzero Fourier coefficients, using $O(k\log^2(2k))$ independent uniform quantum examples and polynomial total computation and working space in the input dimension and $k$. Each quantum example is measured separately. The algorithm uses no membership queries to the target and no inverse of a state-preparation circuit. Fresh residual labels estimate the coefficients used to weight a subsequent independent batch of Fourier samples. A majority of signs of the resulting averages is an explicitly evaluable Boolean hypothesis. Classical heavy-Fourier search on this known hypothesis produces candidate indices, which a separate label batch certifies against the original target. Fresh certification protects the archive even when a proposal is inaccurate. A capped adaptive progress argument requires only a constant probability of contraction per trial. Pruning prefixes of random frequency hashes and correcting the surviving heavy characters gives classical bit time $\widetilde O(n^2k^3+nk^5+k^7)$, with every parity calculation and exact refit charged. We account for adaptive refitting, exact rational arithmetic, unsuccessful trials and all resource costs. This complements the smaller sample bound in the cited version of Ye by supplying a fixed polynomial computational guarantee. The bounds are asymptotic and do not assert practical running time.
\end{abstract}

\section{Result, access model and contribution}
Let $f:\F_2^n\to\{-1,1\}$ and write
\[
 \chi_s(x)=(-1)^{s\cdot x},\qquad
 c_s=\widehat f(s)=2^{-n}\sum_x f(x)\chi_s(x),\qquad
 S=\supp\widehat f.
\]
All norms and expectations on the cube use its uniform probability measure.
The supplied quantum examples are independent copies of
\begin{equation}\label{eq:example}
 |\psi_f\rangle=2^{-n/2}\sum_x |x\rangle|b_f(x)\rangle,
 \qquad b_f(x)=(1-f(x))/2.
\end{equation}
The input consists only of these copies and integers $n,k$ with $|S|\le k$.
The cost of preparing the supplied examples is outside this access model; no preparation circuit or its inverse is supplied.
A label batch means computational-basis measurements of fresh copies, giving independent uniform pairs $(X,f(X))$.

\begin{theorem}\label{thm:main}
For $n\ge1$ and $1\le k\le2^n$, every Boolean $f$ with $|S|\le k$ can be learned exactly with probability at least $2/3$ using
\[
 O\bigl(k\log^2(2k)\bigr)
\]
supplied uniform quantum examples. The output is the complete list of pairs $(s,c_s)$ with $c_s\ne0$, with exact rational coefficients. Total quantum work is $O(nk\log^2(2k))$. Conservative classical bit-time and live-space bounds are
\begin{equation}\label{eq:resources}
 \widetilde O(n^2k^3+nk^5+k^7)\quad\text{and}\quad
 \widetilde O(nk^2+k^3),
\end{equation}
respectively. Here only fixed powers of logarithms in $n$ and $k$ are suppressed. These resource caps hold on every measurement and random tape; an unsuccessful execution may return $\FAIL$. No chosen-input query to $f$ is used.
\end{theorem}

The theorem is about exact representation recovery, including coefficients that are rarely observed by Fourier sampling. It imposes no degree, Fourier-flatness or variance promise. The case $n=0$ needs one label; an upper bound $k>2^n$ can be replaced by $2^n$.

\paragraph{Prior work and comparison.}
Arunachalam, Chakraborty, Lee, Paraashar and de Wolf~\cite{ACLPW} obtained an $O(k^{3/2}\log^2 k)$ quantum sample bound and a lower bound of $\Omega(k\log k)$ in their stated parameter range. Ye~\cite[version~1]{Ye} improves the upper bound to $O(k\log^2 k/\log\log k)$ by constructing a grouped-sample candidate dictionary and proving support containment. That sample bound is smaller than ours. Ye's Theorem~1.1 states a sample guarantee; it does not state a fixed polynomial computation bound. Its Algorithms~1--2 explicitly construct and solve over the candidate dictionary. We do not assert a computational lower bound for that method. Caro et al.~\cite[Corollary~10]{Caro} give an efficient exact learner with $O(k^4\log(k^2/\delta))$ copies and success $1-\delta$. Thus neither polynomial-time learning alone nor near-linear sample complexity alone is the contribution claimed here.

Our contribution is their simultaneous attainment through a computed approximation with error proportional to the current missing energy. We use Fourier sampling, Boolean granularity~\cite{GOSSW}, concentration inequalities~\cite{Tropp}, and classical Fourier search~\cite{KM} with hashed isolation and local character correction~\cite{Datta} as attributed ingredients. The Fourier search is applied to a stored hypothesis, for which chosen-input evaluations are available. The target itself remains accessible only through supplied examples. We make no claim of optimal logarithmic factors or exhaustive publication priority.

\paragraph{Mechanism.}
The maintained approximation $a$ has a certified support contained in $S$. Fresh labels estimate $c_s$ at an error scale proportional to $\norm{f-a}_2$. These estimates are fixed before drawing new Fourier samples. Thresholded reciprocals of the estimates then define an evaluable, bounded sign-majority hypothesis $H$. Searching $H$ supplies candidate frequencies, including frequencies not directly sampled. Original-target certification and exact refitting protect the archive for every bounded proposal list; accurate proposals additionally contract the missing energy. Separating constant-confidence proposals from high-confidence archive safety permits a capped sequence of attempts, without requiring every proposal to succeed. Keeping $f-a$, rather than $f-H$, as the residual is essential: $H$ can have a dense Fourier spectrum.

\paragraph{Version history.}
Version~0.1.0 used $O(k\log^3(2k))$ copies. Version~0.2.0 separated constant-confidence proposals from archive safety, reducing this to $O(k\log^2(2k))$. Version~0.3.0 shared known-hypothesis evaluations across frozen depth tests, giving $\widetilde O(n^2k^6+nk^7)$ classical bit time. Version~0.4.0 prunes random hash prefixes and locally corrects heavy characters, giving~\eqref{eq:resources}. Exactness, success, example count and live space are retained. The refinement does not assert a new Fourier-search principle or practical running time.

\section{Basic access and certificate lemmas}
Parseval gives $\sum_s c_s^2=1$. We use the following consequence of Boolean granularity~\cite[Theorem~8.1]{GOSSW}:
\begin{equation}\label{eq:granularity}
 c_s\ne0\quad\Longrightarrow\quad |c_s|\ge1/k.
\end{equation}
For actual sparsity $s_0>1$, Theorem~8.1 of the cited April 27, 2011 author manuscript makes every coefficient an integer multiple of $2^{1-\lfloor\log_2 s_0\rfloor}$. A nonzero coefficient therefore has magnitude at least $2/s_0\ge1/k$, proving~\eqref{eq:granularity}. A sparsity-one Boolean function is a signed parity.

Apply Hadamard gates to all registers in~\eqref{eq:example}. The amplitude of outcome $(s,1)$ is $c_s/\sqrt2$. Thus the last bit is one with probability $1/2$, and conditional on that flag, $s$ has probability $c_s^2$. Independent copies give independent flags and conditional samples. In particular, conditioning on enough flags in a fixed batch does not bias the successful indices.

For an archive $U\subseteq S$, put
\begin{equation}\label{eq:energy}
 W(U)=\sum_{s\in S\setminus U}c_s^2,\qquad
 P_Uf=\sum_{s\in U}c_s\chi_s,\qquad \alpha^2=82/81.
\end{equation}
The algorithm does not compute $W(U)$ or $P_Uf$.

\begin{lemma}[Fresh residual estimates]\label{lem:estimates}
Let $a$ be fixed before a fresh label batch, supported on at most $k$ frequencies, and put $r=f-a$, $e=\norm r_2$. Let $\mathcal D$ be any set of at most $P$ frequencies fixed before that batch. For $0<\varepsilon\le1$, $0<\beta<1$, take
\begin{equation}\label{eq:qestimate}
 q\ge32k\varepsilon^{-2}\ln\frac{4(P+1)}\beta.
\end{equation}
Define $v=q^{-1}\sum_\ell r(X_\ell)^2$ and $z_s=q^{-1}\sum_\ell r(X_\ell)\chi_s(X_\ell)$.
With probability at least $1-\beta$,
\begin{equation}\label{eq:estimates}
 e^2/2\le v\le3e^2/2,\qquad
 |z_s-\widehat r(s)|\le\varepsilon e/\sqrt k\quad(s\in\mathcal D).
\end{equation}
This conclusion also holds conditional on any history fixing $a,\mathcal D$ before the batch.
\end{lemma}
\begin{proof}
The residual has at most $2k$ Fourier terms, so $\norm r_\infty\le\sqrt{2k}e\le\sqrt{3k}e$. Each $r(X)\chi_s(X)$ has variance at most $e^2$. Since $|\widehat r(s)|\le e$, its centered absolute bound is at most $(\sqrt{2k}+1)e\le3\sqrt k e$. For the scalar bounded-case Bernstein inequality~\cite[Theorem~6.1(ii)]{Tropp}, at deviation $t=\varepsilon e/\sqrt k$ the denominator satisfies
\[
 \sigma^2+Bt/3\le e^2(1+\varepsilon)\le2e^2.
\]
For $q$ independent centered scalar variables bounded in absolute value by $B$, each with variance at most $\sigma^2$, that inequality bounds the deviation of their average by $2\exp[-qt^2/(2(\sigma^2+Bt/3))]$. Thus a two-sided coefficient failure is at most $2\exp(-q\varepsilon^2/(4k))$.
Also $\operatorname{Var}(r(X)^2)\le3ke^4$ and $|r(X)^2-e^2|\le3ke^2$. At deviation $t=e^2/2$, the denominator is at most $2(3ke^4+ke^4/2)=7ke^4$, giving failure at most $2\exp(-q/(28k))$. Put $L_0=\ln(4(P+1)/\beta)$. Equation~\eqref{eq:qestimate} gives $q\varepsilon^2/(4k)\ge8L_0$ and $q/(28k)\ge L_0$, so the union bound is at most $2(P+1)e^{-L_0}=\beta/2$. If $e=0$, all residual values vanish and the conclusion is exact.
\end{proof}

\begin{lemma}[Original-coefficient certification]\label{lem:certificate}
Suppose $U\subseteq S$, $a$ is supported on $U$, and a candidate list $L$, $|L|\le P$, precedes a fresh batch satisfying~\eqref{eq:qestimate}. Admit $s\in L\setminus U$ precisely when
\begin{equation}\label{eq:certificate}
 z_s^2>4\varepsilon^2v/k.
\end{equation}
On~\eqref{eq:estimates}, no frequency outside $S$ is admitted, and every proposed frequency with $|c_s|>4\varepsilon e/\sqrt k$ is admitted.
\end{lemma}
\begin{proof}
Outside $U$ the mean of $z_s$ is $c_s$. For $c_s=0$, its squared estimate is at most $\varepsilon^2e^2/k$, while the threshold is at least $2\varepsilon^2e^2/k$. If $|c_s|>4\varepsilon e/\sqrt k$, the estimated magnitude exceeds $3\varepsilon e/\sqrt k$, and the square root of the threshold is at most $\sqrt6\varepsilon e/\sqrt k$. The zero-residual case admits nothing. More generally, in any batch $v=0$ forces every observed residual and every $z_s$ to vanish, so no index is admitted even without the accuracy event.
\end{proof}

\begin{lemma}[One Gram event controls adaptive refitting]\label{lem:fit}
For a fresh batch of
\begin{equation}\label{eq:qgram}
 q\ge300(k+1)\ln\frac{2(k+1)}\beta
\end{equation}
uniform labels, with probability at least $1-\beta$, every subset $U\subseteq S$ has a nonsingular empirical Gram matrix, and its exact least-squares fit $a$ satisfies
\begin{equation}\label{eq:invariant}
 \norm{P_Uf-a}_2\le\frac{\sqrt{W(U)}}9,\qquad
 \norm{f-a}_2^2\le\alpha^2 W(U),\qquad \norm a_2^2\le\alpha^2<2.
\end{equation}
The set $U$ may be chosen using this same batch. The empty fit is zero.
\end{lemma}
\begin{proof}
Let $S_0=S\cup\{0\}$ and $p=|S_0|\le k+1$. The analysis-only row $\phi(X)=(\chi_s(X))_{s\in S_0}$ satisfies $\norm{\phi(X)}_2^2=p$ and $\E\phi(X)\phi(X)^\mathsf T=I_p$. The independent positive semidefinite summands $\phi(X_\ell)\phi(X_\ell)^\mathsf T$ have largest eigenvalue $p$ and their expected sum is $qI_p$. Matrix Chernoff~\cite[Corollary~5.2]{Tropp} therefore bounds the lower and upper eigenvalue tails by
\[
 p\left[\frac{e^{-\delta}}{(1-\delta)^{1-\delta}}\right]^{q/p},
 \qquad
 p\left[\frac{e^{\delta}}{(1+\delta)^{1+\delta}}\right]^{q/p},
 \quad 0<\delta<1,
\]
respectively. The scalar exponents satisfy
\[
 \delta+(1-\delta)\ln(1-\delta)\ge\delta^2/2,\qquad
 (1+\delta)\ln(1+\delta)-\delta\ge\delta^2/3.
\]
The first follows by integrating $-\ln(1-t)\ge t$; the second by integrating $\ln(1+t)\ge t-t^2/2$ on $[0,\delta]$. Summing the tails gives
\[
 \Pr\!\left[\norm{q^{-1}A_{S_0}^{\mathsf T}A_{S_0}-I_p}_{\rm op}>\delta\right]
 \le2p\exp\!\left(-\frac{q\delta^2}{3p}\right),\quad0<\delta<1.
\]
Take $\delta=1/10$; the budget in~\eqref{eq:qgram} is at least $300p\ln(2p/\beta)$ because both positive factors increase with $p$. This single event controls all principal submatrices and cross blocks, without a union bound over subsets. For $V=S\setminus U$, exact pointwise Fourier inversion gives the label vector $y=A_Uc_U+A_Vc_V$. The normal equations therefore give
\[
 a_U-c_U=(A_U^{\mathsf T}A_U/q)^{-1}(A_U^{\mathsf T}A_V/q)c_V.
\]
The two operator norms are at most $10/9$ and $1/10$, respectively. This proves the first inequality; orthogonality proves the second. Finally
\[
 \norm a_2\le\sqrt{1-W(U)}+\frac{\sqrt{W(U)}}9\le\sqrt{1+1/81}.
\]
The proof uses unknown $S_0$ only to define an event. The algorithm constructs only the columns it actually retains. If $U=S$, the fit is exactly $f$.
\end{proof}

\section{Fourier search on a known Boolean function}
Prefix-energy search is associated with Kushilevitz--Mansour~\cite{KM}; hashed isolation and local character correction appear in Datta et al.~\cite[Lemma~8 and Claims~12--15]{Datta}. We combine pruned hash-prefix tests with explicit coordinate recovery and account for all processing. The proof below is self-contained. Only genuinely heavy frequencies need isolation: extra proposals are permitted and will face original-target certification. The known Boolean hypothesis may have a dense Fourier spectrum.

\begin{lemma}[Pruned hashed heavy-spectrum recovery]\label{lem:search}
Given deterministic evaluation access to $H:\F_2^n\to\{-1,1\}$, $n\ge1$, rational $0<\tau\le1$ and $0<\eta<1$, set
\begin{align}\label{eq:hashparameters}
 P&=\lceil4/\tau^2\rceil,&
 t&=\left\lceil\log_2\frac{4096}{\eta\tau^4}\right\rceil,\notag\\
 L_0&=\left\lceil\log_2\frac{16tP}{\eta}\right\rceil,&
 N&=\lceil32\tau^{-4}L_0\rceil,\qquad M=\lceil224\tau^{-2}\rceil.
\end{align}
Let $J$ be the least positive odd integer with $2^J\ge4nP/\eta$.
A classical algorithm returns at most $P$ distinct explicit $n$-bit indices, or $\FAIL$. Except with probability at most $3\eta/4$, it does not fail and its list contains every $s$ with $|\widehat H(s)|\ge\tau$. It uses at most $2tN+2nJM$ evaluations of $H$. All caps hold on every tape, and the success guarantee holds conditional on any preceding history fixing $H$.
\end{lemma}
\begin{proof}
Write $d_s=\widehat H(s)$, so $\sum_s d_s^2=1$. Sample $t$ independent uniform rows $b_1,\ldots,b_t\in\F_2^n$ to form a binary matrix $B$. Let $B_j$ consist of its first $j$ rows. Do not condition on full rank or discard dependent rows, including when $t>n$.

\paragraph{Bucket and projection identities.}
For labels $v$ of the indicated lengths define
\begin{align}\label{eq:bucket}
 E_{j,v}&=\sum_{s:B_js=v}d_s^2
 =\E_{x,u}H(x)H(x+B_j^{\mathsf T}u)(-1)^{v\cdot u},\notag\\
 \mathcal P_v(y)&=\sum_{s:Bs=v}d_s\chi_s(y)
 =\E_u H(y+B^{\mathsf T}u)(-1)^{v\cdot u}.
\end{align}
In the first identity $x\in\F_2^n,u\in\F_2^j$ are independent uniform points; in the second $u\in\F_2^t$ is uniform. Expand the Fourier series. Orthogonality over $x$ removes the first identity's cross terms, and averaging $(-1)^{u\cdot(B_js+v)}$ gives the indicator $B_js=v$. The same averaging proves the second identity. Uniformity is in the full label-parameter space, so the identities also hold for dependent rows and infeasible labels. No inverse matrix or coset representative is needed.

\paragraph{Prune frozen hash prefixes.}
Start with the empty label. At depth $j=1,\ldots,t$, freeze both children of every retained label into $C_j$, with $|C_j|\le2P$. If empty, return the empty list. Initialize integer counters $A_v=0$. For each of $N$ fresh independent uniform pairs $(x,u)\in\F_2^n\times\F_2^j$, compute $B_j^{\mathsf T}u$ and $H(x)H(x+B_j^{\mathsf T}u)$ once. Add this product times $(-1)^{v\cdot u}$ to every frozen label's counter; then discard the row. After the full batch retain labels with $A_v/N\ge3\tau^2/4$. Return $\FAIL$ if their number exceeds $P$. Otherwise choose the next children only after this batch is complete. Write $\mathcal C$ for the final distinct $t$-bit labels.

Conditional on $B$ and the preceding history, each tested label's independent row summands lie in $[-1,1]$ and have mean $E_{j,v}$. The bounded-variable Hoeffding estimate at deviation $\tau^2/4$ gives
\[
 \Pr[|A_v/N-E_{j,v}|>\tau^2/4\mid\text{past},B]
 \le2\exp(-N\tau^4/32)\le\frac{\eta}{8tP}.
\]
Here $L_0\ge\ln(16tP/\eta)$. Union-bound over the frozen labels, then take conditional expectations and sum over the $t$ levels; failures at unreached levels are defined false. The probability of any reached inaccurate estimate is at most $\eta/4$, for every fixed $B$. No independence among labels or conditioning on future accuracy is used.

On simultaneous accuracy, every prefix of a coefficient with magnitude at least $\tau$ has energy at least $\tau^2$ and survives. Every retained prefix has energy at least $\tau^2/2$. Same-depth buckets are disjoint with total energy one, so at most $2/\tau^2\le P$ survive. Thus no overflow occurs, and all heavy coefficients' full labels belong to $\mathcal C$. Only $2P$ labels per depth are ever tested; there is no iteration over $2^t$ labels.

\paragraph{Isolate only heavy coefficients.}
For each fixed heavy $s$, define the noise energy in its bucket by
\[
 T_s(B)=\sum_{a\ne s:\,Ba=Bs}d_a^2.
\]
For $a\ne s$, independent uniform rows give $\Pr[Ba=Bs]=2^{-t}$, regardless of $n$ versus $t$. Consequently $\E_B T_s\le2^{-t}$. There are at most $\tau^{-2}$ heavy coefficients. Markov and a union bound imply
\begin{equation}\label{eq:hashisolation}
 \Pr_B[\exists s:\ |d_s|\ge\tau,\ T_s>\tau^2/1024]
 \le\frac{1024}{2^t\tau^4}\le\frac\eta4.
\end{equation}
On the complementary good-hash event, heavy coefficients have distinct full labels: a collision with another heavy $s'$ would give $T_s\ge d_{s'}^2\ge\tau^2$. Buckets without a heavy coefficient require no isolation guarantee and may yield arbitrary proposals.

\paragraph{Recover original coordinate bits.}
For each coordinate $i=1,\ldots,n$, perform $J$ repetitions. In each, choose a fresh uniform $y\in\F_2^n$ and $M$ independent uniform vectors $u\in\F_2^t$. Using~\eqref{eq:bucket}, accumulate both projections $\mathcal P_v(y)$ and $\mathcal P_v(y+e_i)$ for all $v\in\mathcal C$, sharing each row's two $H$ evaluations. For each label, the vote is the product of the signs of its two integer sums, with $\sgn(0)=1$. A positive majority of the odd $J$ votes sets bit $i$ of the decoded index $s_v$ to zero; a negative majority sets it to one. All repetitions and coordinates use fresh randomness. Keep only the current coordinate's vote counters and the bits already decoded. Finally discard $s_v$ if $Bs_v\ne v$, and return the remaining indices. Different retained labels cannot then produce the same index.

Fix a good $B$, a completed pruning transcript and a surviving heavy label $v=Bs$. Write
\[
 \mathcal P_v(y)=d_s\chi_s(y)+r_s(y),\qquad
 \norm{r_s}_2^2=T_s\le\tau^2/1024.
\]
The residual need not have mean zero. Since both $y$ and $y+e_i$ are uniform, Markov and a union bound give
\[
 \Pr_y[|r_s(y)|>\tau/4\ \text{or}\ |r_s(y+e_i)|>\tau/4]
 \le32T_s/\tau^2\le1/32.
\]
Conditional on any $y$, each empirical projection errs by more than $\tau/4$ with probability at most
\[
 2\exp(-M\tau^2/32)\le2e^{-7}\le1/64.
\]
The two estimates may be correlated; their failure union is at most $1/32$. Outside these events, each estimated projection is within $\tau/2$ of its $d_s\chi_s$ term. As $|d_s|\ge\tau$, both signs are correct and their product equals
\[
 \sgn(d_s)^2\chi_s(y)\chi_s(y+e_i)
 =\chi_s(e_i)=(-1)^{s_i}.
\]
This includes negative coefficients and the constant character. A vote is wrong with probability at most $1/16$. For a fixed heavy index and coordinate, the $J$ repetitions are independent conditional on $H,B$ and the pruning transcript. A union bound over erroneous-majority subsets gives
\[
 \Pr[\text{decoded bit wrong}]\le2^J(1/16)^{\lceil J/2\rceil}
 \le2^{-J}\le\frac{\eta}{4nP}.
\]
Union-bound over at most $nP$ heavy bits. No independence across coordinates or labels is required. Except with probability $\eta/4$, every heavy index is recovered exactly and passes $Bs_v=v$.

The pruning, hash and conditional correction budgets sum to $3\eta/4$. Bad events are charged without any correctness assertion on those events. All loops, lists and counters remain capped there. Each pruning row uses two evaluations and each correction row two, giving $2tN+2nJM$. The entire argument holds for every fixed $H$, hence also conditional on earlier histories fixing it before the new random bits. All function queries are to the known $H$; no target examples are consumed.
\end{proof}

\paragraph{Complete processing and storage.}
Let $C_H$ cap evaluation bit time, $S_H$ include the stored evaluator and its workspace, $d_0$ bound the rational parameter lengths, and $h_s=O(d_0+t+\log(2nNJMP))$ bound counter and loop widths. Forming $B_j^{\mathsf T}u$ costs $O(nj)$ bit operations by scanning rows. All prefix-label parities, integer increments, exact threshold comparisons and prefix copying give the conservative cost
\[
 O\bigl(nt+tN(C_H+nt+P(t+h_s))+tP(d_0+h_s)^2+Pt^2\bigr).
\]
Correction pays every matrix-vector product, input XOR and bucket parity; vote handling and final hash checks give
\[
 O\bigl(nJM(C_H+nt+P(t+h_s))+nJP h_s+ntP\bigr).
\]
Parameter construction uses integer comparisons with powers of two and has polynomial cost in $d_0+\log n$. Random-bit consumption is at most
\begin{equation}\label{eq:hashrandom}
 nt+N\bigl[tn+t(t+1)/2\bigr]+nJ(n+Mt).
\end{equation}
Store the $nt$ hash bits, current labels/counters, decoded $n$-bit indices, input buffers and one coordinate's votes. The live-space bound is
\[
 S_H+O\bigl(nt+P(n+t+h_s)+n+t+h_s\bigr)
   +\operatorname{poly}(d_0+\log n).
\]
Projection counters are integers of magnitude at most $M$; only their signs are used. No growing rational projection, all-label array or sample table is stored. Decoded-bit storage scans fit the stated correction bound because $M\ge P$. The specialized parameters and total learner costs are substituted in Section~\ref{sec:costs}.

\section{Constructing a queryable majority hypothesis}
Assume $U\subseteq S$ and a fixed, known rational fit $a$ satisfies~\eqref{eq:invariant}. Write $W=W(U)$ and $e=\norm{f-a}_2$. Draw an estimator label batch with Lemma~\ref{lem:estimates}, using list-size argument $k$ in that lemma. For any index $s$, define
\begin{equation}\label{eq:ctilde}
 \widetilde c_s=\widehat a(s)+q_e^{-1}\sum_\ell
 (f(X_\ell)-a(X_\ell))\chi_s(X_\ell).
\end{equation}
The formula is evaluated only at indices subsequently obtained from fresh Fourier measurements. Its values at all other indices need never be computed.

Conditional on the entire earlier history, $S$ is fixed. Apply Lemma~\ref{lem:estimates} to $\mathcal D=S$ in the analysis. This yields the event
\begin{equation}\label{eq:uniformtrue}
 e^2/2\le v\le3e^2/2,\qquad
 |\widetilde c_s-c_s|\le\Delta:=\varepsilon e/\sqrt k\quad(s\in S),
\end{equation}
with probability at least $1-\beta$. The union bound does not require knowing $S$. It is not a claim of uniform estimation over all $2^n$ indices.

After fixing the estimator batch, draw $g$ independent groups of $m$ Fourier samples $s_{ij}$, retaining multiplicities. A frequency passes the cutoff if $\widetilde c_s\ne0$ and
\begin{equation}\label{eq:cutoff}
 \widetilde c_s^2\ge16\varepsilon^2v/k.
\end{equation}
Define $w_s=1/\widetilde c_s$ on a pass and $w_s=0$ otherwise. Form
\begin{equation}\label{eq:H}
 Q_i(x)=\frac1m\sum_{j=1}^m w_{s_{ij}}\chi_{s_{ij}}(x),\qquad
 H(x)=\operatorname{majority}_{1\le i\le g}\sgn Q_i(x),
\end{equation}
where $g$ is odd and $\sgn(0)=1$. Zero estimates are never divided. This is an explicit finite classical circuit with exact rational weights. It has no call to $f$.

\begin{lemma}[Bias and variance of the computed averages]\label{lem:bias}
Suppose $W>0$ and~\eqref{eq:uniformtrue} holds. Conditional on the estimator batch, let $h=\E Q_i$ as a function on the cube. Then
\begin{equation}\label{eq:bias}
 \norm{h-f}_2^2\le36\varepsilon^2e^2,\qquad
 \operatorname{Var}(Q_i(x))\le\frac{9k}{4m}\quad\text{for every }x.
\end{equation}
\end{lemma}
\begin{proof}
The cutoff magnitude $t=4\varepsilon\sqrt{v/k}$ obeys
$2\sqrt2\Delta\le t\le2\sqrt6\Delta<5\Delta$.
On a pass, $|c_s/\widetilde c_s|\le1+\Delta/t<3/2$ and the signs agree. On a failure, $|c_s|<6\Delta$, including when the estimate is zero. Since future Fourier samples have probabilities $c_s^2$,
\[
 h(x)=\sum_{s\in S:\,s\text{ passes}}\frac{c_s^2}{\widetilde c_s}\chi_s(x).
\]
The coefficient error on a pass is $c_s(c_s-\widetilde c_s)/\widetilde c_s$, of magnitude at most $3\Delta/2$. On a failure it is at most $6\Delta$. The passing and failing sets partition $S$, so Parseval bounds the total squared error by $36|S|\Delta^2\le36\varepsilon^2e^2$. A single weighted sample has second moment
$\sum_{s\text{ passes}}c_s^2/\widetilde c_s^2\le9k/4$, independently of $x$. Independence within each group proves the second inequality. The function $h$ is analysis-only; the algorithm evaluates $H$, not $h$.
\end{proof}

\begin{lemma}[Bounded majority and original residual scale]\label{lem:majority}
Take $m=144k$. Conditional on~\eqref{eq:uniformtrue},
\begin{equation}\label{eq:majorityexpect}
 \E_{\rm groups}\norm{H-f}_2^2\le576\varepsilon^2e^2+4\cdot2^{-g}.
\end{equation}
If $\varepsilon^2\le\beta/262144$ and $2^g\ge512k^2/\beta$, then for $W>0$,
\begin{equation}\label{eq:Hgood}
 \Pr_{\rm groups}\!\left[\norm{H-f}_2^2>W/64\right]\le\beta.
\end{equation}
\end{lemma}
\begin{proof}
Let $A=\{x:|h(x)-f(x)|>1/2\}$. Its measure is at most $4\norm{h-f}_2^2$. Outside $A$, $h$ has the correct sign with margin at least $1/2$. Lemma~\ref{lem:bias} and Chebyshev bound each group's sign-error probability there by $1/16$. For a fixed such $x$, the groups are independent conditional on the estimator batch. Summing over possible erroneous majorities bounds their probability by
$2^g(1/16)^{g/2}=2^{-g}$.
No independence over different inputs $x$ is asserted. Integrating, and using the squared error bound four for Boolean outputs, gives
\[
 \E_{\rm groups}\norm{H-f}_2^2
 \le16\norm{h-f}_2^2+4\cdot2^{-g},
\]
which proves~\eqref{eq:majorityexpect}. Since $e^2\le\alpha^2 W$ and $\alpha^2<2$, the first term is at most $\beta W/128$. By~\eqref{eq:granularity}, positive $W$ is at least $1/k^2$, so the second term is also at most $\beta W/128$. Markov gives~\eqref{eq:Hgood}.
\end{proof}

The majority in~\eqref{eq:H} is taken after taking signs. Bounded output controls the exceptional set $A$ even if the rational averages are large there. This argument does not assume an unbiased computed hypothesis or stability of an unbounded majority polynomial.

\section{Complete algorithm and success probability}
The generic lemmas above allow their failure parameters to be chosen separately. Use constant proposal confidence $\rho$ and smaller archive-safety failure budget $\beta$:
\begin{align}\label{eq:parameters}
 \ell&=\lceil\log_2(2k)\rceil,&T&=8\ell,&\beta&=\frac1{200T},\notag\\
 \rho&=\frac1{16},&\varepsilon&=\frac1{2048},&\gamma&=\frac1{32},\notag\\
 m&=144k,&\tau&=\frac1{4k},&P&=64k^2.
\end{align}
Choose the least positive odd $g$ such that $2^g\ge512k^2/\rho=8192k^2$, and set
\begin{align}\label{eq:batches}
 q_e&=\left\lceil32k\varepsilon^{-2}\ln\frac{4(k+1)}\rho\right\rceil,\notag\\
 q_c&=\max\left\{
 \left\lceil32k\gamma^{-2}\ln\frac{4(P+1)}\beta\right\rceil,
 \left\lceil300(k+1)\ln\frac{2(k+1)}\beta\right\rceil\right\}.
\end{align}
For an exact finite implementation, replace each $\ln z$ in a sample count by $\lceil\log_2 z\rceil$. This only increases the count and is computed by integer comparisons. The same convention applies to search and final validation. In Lemma~\ref{lem:estimates}, the proposal uses list-size bound $k$ and failure budget $\rho$, while certification uses $P$, accuracy $\gamma$ and failure budget $\beta$. The majority lemma's failure parameter and the search lemma's $\eta$ are instantiated as $\rho$.

Initialize $U=\varnothing$ and $a=0$. Execute exactly $T$ trials unless an explicitly stated global abort returns $\FAIL$:
\begin{enumerate}
\item Measure $q_e$ fresh labels. Store the residuals and $v$ and define~\eqref{eq:ctilde}. These estimates precede the new Fourier samples.
\item Fourier-measure $4gm$ new copies, discarding every flag-zero outcome. If fewer than $gm$ flags succeed, set $L=\varnothing$ and proceed to step~4. Otherwise split the first $gm$ successful indices into $g$ groups. If more than $k$ distinct grouped indices occur, again set $L=\varnothing$ and proceed to step~4. This distinct-count guard cannot fire for ideal examples of an admitted target. Compute the estimates at these indices and construct~\eqref{eq:H}.
\item Apply Lemma~\ref{lem:search} to this known $H$ with $\tau,P$ and failure budget $\rho$. Use its returned list, or set $L=\varnothing$ if the search returns $\FAIL$. A failed proposal does not abort the learner or modify its archive.
\item Measure $q_c$ new labels, including when $L$ is empty. Apply~\eqref{eq:certificate} to $L\setminus U$ with parameter $\gamma$, using the old residual $f-a$. Add all admissions permanently, including indices whose fitted coefficients may later be zero. Globally abort if $|U|>k$. Refit the original target labels $f(X_j)$ from this same batch exactly on the updated archive; globally abort on a singular nonempty system or $\norm a_2^2>2$. The empty fit is zero.
\end{enumerate}
After the last trial, remove zero fitted coefficients, check sparsity at most $k$, and verify $a^2=1$ by exact Fourier convolution. The convolution has at most $k^2$ pairs; it does not enumerate the cube. Finally compare $a$ with $\lceil2k\ln100\rceil$ fresh labels. Return the exact list if all checks pass, and otherwise return $\FAIL$. No stopping decision uses the unknown energy $W$ or an inference from sampled $v=0$.

\begin{proposition}[Safety for arbitrary proposals]\label{prop:safety}
Condition on any history with $U\subseteq S$ and~\eqref{eq:invariant}. For every proposal list $L$ of size at most $P$ fixed before step~4, let $C$ be the residual accuracy event of Lemma~\ref{lem:estimates} with parameter $\gamma$, and let $G$ be the full-support Gram event of Lemma~\ref{lem:fit}. Their combined failure probability is at most $2\beta$. On $C\cap G$, no global abort occurs, the updated state satisfies~\eqref{eq:invariant}, and
\[
 U\subseteq U^+\subseteq S,\qquad W(U^+)\le W(U).
\]
If $W(U)=0$, the new fit remains exactly $f$.
\end{proposition}
\begin{proof}
The list precedes the fresh certification batch, even if it depends on all earlier data and an inaccurate $H$. On $C$, Lemma~\ref{lem:certificate} admits only true indices. Thus $|U^+|\le k$. On $G$, the nonsingularity, error and norm guarantees hold for every subset of $S$, including $U^+$ chosen using these same rows. The fit therefore satisfies~\eqref{eq:invariant} and its squared norm is less than two. No archive index is deleted, so missing original energy cannot increase. When $W=0$, the old fit is $f$, all certificate residuals are identically zero and no new index is admitted; the Gram event refits $f$ exactly. The proposal may still contain indices outside $S$; zero certificate residuals reject these as well. The events $C,G$ need not be independent. Their probabilities are bounded conditionally on each history preceding their shared fresh batch, and a union bound gives $2\beta$.
\end{proof}

\begin{proposition}[Constant conditional progress]\label{prop:contraction}
For every history satisfying the preceding invariant with $W(U)>0$, a trial safely preserves that invariant and gives $W(U^+)<W(U)/16$ with probability at least $3/4-2\beta\ge1/2$.
\end{proposition}
\begin{proof}
Use failure budget $\rho$ in the proposal estimates, majority lemma and hashed search (Lemma~\ref{lem:search}, failure at most $3\rho/4$). Since $\varepsilon^2=\rho/262144$, Lemma~\ref{lem:majority} gives $\norm{H-f}_2^2\le W/64$ except with conditional probability at most $\rho$ on the estimator event. The flag count is below $gm$ with probability at most $\exp(-gm/4)\le\rho$; conditioning on enough flags preserves the independent Fourier law of the successful indices. A staged union bound over estimates, flags, majority accuracy and search gives simultaneous success with probability at least $1-4\rho=3/4$. No independence among these events is assumed.

On those four events, a true coefficient omitted from $L$ has $|\widehat H(s)|<1/(4k)\le|c_s|/4$, and hence
\begin{equation}\label{eq:omitted}
 \sum_{s\in S\setminus L}c_s^2
 \le\frac{16}{9}\norm{H-f}_2^2\le W/36.
\end{equation}
On the fresh certification event $C$, every proposed true frequency still not admitted has magnitude at most $4\gamma e/\sqrt k$. Their total energy is at most $16\gamma^2e^2\le\alpha^2W/64$. Therefore
\begin{equation}\label{eq:contract}
 W(U^+)\le\left(\frac1{36}+\frac{\alpha^2}{64}\right)W<W/16.
\end{equation}
Intersecting also with $G$ gives a safe contraction. The probability is at least $3/4-2\beta\ge1/2$, because $T\ge8$ and $\beta\le1/1600$. All other safety-preserving trials remain nonincreasing in $W$, whether or not the proposal was accurate.
\end{proof}

\begin{proof}[Proof of the success claim in Theorem~\ref{thm:main}]
Initially $W=1$ and~\eqref{eq:invariant} holds. Let $B$ be the event that a first safety failure $C^c$ or $G^c$ occurs during the $T$ trials. Before the first such failure the state is valid, regardless of earlier proposal quality. Proposition~\ref{prop:safety} and a first-failure union bound give
\begin{equation}\label{eq:safetyfailure}
 \Pr(B)\le2T\beta=1/100.
\end{equation}
We do not condition the entire process on future safety before counting progress.

For the proof only, define indicators $I_1,\ldots,I_T$. If a safety failure occurred earlier, or $W$ is already zero at the trial's start, set $I_t=1$. Otherwise set $I_t=1$ precisely when the four good proposal events from Proposition~\ref{prop:contraction} and $C\cap G$ occur. If execution terminates, unused indicators are one. A global abort cannot precede a safety failure. Thus, for the filtration $\mathcal F_{t-1}$ generated by the past trials, either $I_t=1$ by convention or the current state is valid with positive $W$, and
\[
 \Pr(I_t=1\mid\mathcal F_{t-1})\ge1/2.
\]
Put $Z_t=\sum_{j\le t}I_j$. Conditional expectation yields
\[
 \E[2^{-I_t}\mid\mathcal F_{t-1}]
 =1-\tfrac12\Pr(I_t=1\mid\mathcal F_{t-1})\le3/4,
 \qquad \E[2^{-Z_T}]\le(3/4)^T.
\]
This iteration does not assume independent trials. By Markov,
\begin{equation}\label{eq:progresscount}
 \Pr(Z_T<\ell)\le2^\ell(3/4)^T
 =\bigl[2(3/4)^8\bigr]^\ell<1/4,
\end{equation}
since $\ell\ge1$ and $2(3/4)^8=6561/32768<1/4$.

On $B^c$, zero missing energy is absorbing. If it has not already occurred, $Z_T\ge\ell$ supplies at least $\ell$ genuine factor-16 contractions, interspersed with nonincreasing safe trials. Consequently
\[
 W(U_T)<16^{-\ell}\le\frac1{16k^4}<\frac1{k^2}.
\]
Granularity forces $W(U_T)=0$. The final exact refit is therefore $f$, and every output check passes surely. Combining~\eqref{eq:safetyfailure} and~\eqref{eq:progresscount} gives exact-output probability greater than $1-1/100-1/4=37/50>2/3$. Dummy indicators after a safety failure do not imply correctness on that event; its probability is charged separately.

For completeness, the final validation limits erroneous acceptance. If two distinct Boolean functions have sparsity at most $k$, their difference $d$ has at most $2k$ Fourier terms and $\norm d_\infty=2$. Cauchy--Schwarz gives $4\le2k\norm d_2^2$, so their disagreement probability is at least $1/(2k)$. Conditional on any incorrect candidate passing Booleanity, fresh validation accepts it with probability at most $1/100$. The exact-output success bound above does not depend on this additional error-detection guarantee.
\end{proof}

\section{All resource costs and exact arithmetic}\label{sec:costs}
The proposal parameters $\rho,\varepsilon,\gamma$ are absolute constants. We have $T,\beta^{-1},g=O(\log(2k))$ and $q_e,q_c=O(k\log(2k))$. Including every unsuccessful proposal, refit and final validation, the total supplied-copy cap is
\begin{equation}\label{eq:copies}
 T(q_e+4gm+q_c)+\lceil2k\ln100\rceil=O(k\log^2(2k)).
\end{equation}
The logarithm of $n$ in the search repetition count $J$ does not enter this count: those evaluations are of $H$. Each copy requires at most $O(n)$ single-qubit gates and measurements. No consumed copy is reused and no collective measurement is required.

\paragraph{Refits and coefficient precision.}
An exact refit has at most $k$ sign columns and $q_c$ original sign labels. Its integer Gram matrix and right side have entries bounded by $q_c$. A minor of order at most $k$ has magnitude at most $k!q_c^k$. Thus each coefficient has a common denominator and numerator of bit length
\[
 b_a=O\bigl(k\log(2kq_c)\bigr)=\widetilde O(k).
\]
For an explicit conservative implementation, compute the Gram determinant and at most $k$ column-replacement determinants successively by fraction-free elimination~\cite{Bareiss}, then use Cramer's rule. A zero Gram determinant returns $\FAIL$ for a nonempty fit; the empty fit is zero. Zero column-replacement determinants simply give zero coefficients.

The integer-precision guarantee follows from bordered minors. Let $M$ denote the Gram matrix or a column-replaced matrix being eliminated. If $K$ is the previous invertible pivot block, $\Delta=\det K$ and $D_{uv}$ is the determinant of $K$ bordered by row $u$ and column $v$ of $M$, Schur elimination gives $D_{uv}=\Delta S_{uv}$. Hence
\[
 D_{pp}D_{ij}-D_{ip}D_{pj}
 =\Delta\det M[K\text{ bordered by rows }p,i\text{ and columns }p,j].
\]
This identity makes the fraction-free update's division by the previous pivot exact; the initial empty determinant is one. Row swaps among remaining rows preserve the identity and track the determinant sign. If no nonzero pivot remains in the current column, the determinant is zero. Inductively stored entries are signed minors, so they and temporary products have $O(b_a)$ bits. Schoolbook multiplication and exact division cost $O(b_a^2)$. One determinant takes $O(k^3b_a^2)$ time and $O(k^2b_a)$ bits; at most $k+1$ sequential determinants therefore take $\widetilde O(k^6)$ time and $\widetilde O(k^3)$ space. Norm checking is an exact sum-of-squares comparison. This bound holds on singular and statistically unsuccessful tapes too. Each original-label refit resets precision, so denominators do not compound over trials.

Writing $a_s=A_s/D$ with $D>0$, a row residual has numerator $D f(x)-\sum_s A_s\chi_s(x)$ of $O(b_a+\log(2k))$ bits. Requested estimates have common denominator $q_eD$ and numerators of length $O(b_a+\log(2kq_e))$. Squared residuals and certificate statistics have the same order of bit length, allowing a constant factor for squares. All cutoffs, certification, norms and Booleanity checks use exact integer cross-products.

\paragraph{Evaluating the hypothesis.}
At most $k$ distinct grouped Fourier indices are stored, with group multiplicities. The distinct-count cap cannot fire for ideal examples because flag-one indices belong to $S$; flag-zero outcomes are discarded. For passing nonzero estimates write $\widetilde c_s=N_s/D_e$ with $D_e>0$. The product $R$ of their $|N_s|$ has
\[
 B_w=O\bigl(k[b_a+\log(2kq_e)]\bigr)=\widetilde O(k^2)
\]
bits; the empty product is one. Use integer weight $\sgn(N_s)R/|N_s|$ and the group multiplicity. The common positive factor $D_e/(mR)$ does not affect a sign. Zero estimates and failed cutoffs receive weight zero. Product formation and at most $k$ long divisions cost at most $O(kB_w^2)=\widetilde O(k^5)$ bit operations. Store the resulting group weights once, in $\widetilde O(k^3)$ bits.

One $H$ evaluation computes at most $k$ explicit $n$-bit parities and $g$ sums of at most $k$ $B_w$-bit weighted signs, for $\widetilde O(nk+k^3)$ bit operations. Ties give $+1$. Even inaccurate or tiny estimates yield a deterministic Boolean evaluator with these exact representation bounds; no statistical conditioning is needed for its running time. Earlier hypothesis circuits are not recursively embedded.

\paragraph{Complete classical time and randomness.}
For search, $\tau=1/(4k)$, $P=64k^2$ and $\eta=\rho=1/16$. Its finite parameters are
\begin{align*}
 t&=\lceil\log_2(2^{24}k^4)\rceil,&
 L_0&=\lceil\log_2(16384tk^2)\rceil,\\
 N&=8192k^4L_0,& M&=3584k^2,
\end{align*}
with $J$ the least positive odd integer satisfying $2^J\ge4096nk^2$. Thus $t,L_0=O(\log(2k))$ and $J=O(\log(2nk))$. The $2tN+2nJM$ evaluations cost
\[
 \widetilde O(nk^5+k^7+n^2k^3+nk^5).
\]
Pruning input/hash-vector formation costs $\widetilde O(nk^4)$ and its full label-parity scan $\widetilde O(k^6)$. Correction input/hash-vector formation costs $\widetilde O(n^2k^2)$ and its full label-parity scan $\widetilde O(nk^4)$. Parameter arithmetic, vote handling, decoded-bit storage and exact hash consistency checks fit these bounds. No parity or hash operation is free and no $2^t$-label enumeration is performed.

For clarity, the remaining per-trial work has the following conservative bit costs. Reading and grouping measured outcomes, even by linear index comparisons, costs $\widetilde O(nk^2)$. Evaluating the old fit and residual energy on $q_e$ or $q_c$ rows, and computing at most $k$ requested proposal estimates, costs $\widetilde O(nk^2+k^3)$. Fresh certification of at most $P$ candidates costs $\widetilde O(nk^3+k^4)$ for all parities, accumulations, squared tests and archive checks. Forming the at-most-$k$-column Gram system costs $\widetilde O(nk^2+k^3)$; its exact fit costs $\widetilde O(k^6)$ as above. No $P$-column fit is formed: archive overflow aborts first. The $\widetilde O(k^5)$ inverse-weight construction is also paid.

Final Booleanity uses at most $k^2$ XOR-keyed coefficient products. Forming products and sorting explicit keys/payloads costs $\widetilde O(nk^2+k^4)$; accumulated numerators gain only $O(\log k)$ bits. Fresh final validation and exact output cost $\widetilde O(nk^2+k^3)$. This does not enumerate the cube. Every listed term fits $\widetilde O(n^2k^3+nk^5+k^7)$ for $n,k\ge1$. Multiplication by $T=O(\log(2k))$ preserves that bound and pays all unsuccessful trials. Equation~\eqref{eq:hashrandom} gives $\widetilde O(nk^4+n^2)$ total classical search random bits, including the $T$ trials, generated and discarded row by row.

\paragraph{Live space and unsuccessful tapes.}
Retain only current labels/candidates, label rows and residuals, the archive, group indices/weights, certificate accumulators, hash rows, decoded indices and a reused determinant workspace. Their sum is $\widetilde O(nk^2+k^3)$ bits. The hash uses $O(n\log(2k))$ bits; search/projection/vote counters have logarithmic width, certificate accumulators have $\widetilde O(k)$ bits each, and determinant storage has $\widetilde O(k^3)$ bits. The final $k^2$ keyed products fit the same envelope. Old hypotheses and batches are discarded after refitting. No search sample table or full trial transcript is required.

Every trial, batch, label count and arithmetic input size is capped. A failed proposal yields an empty list and still pays fresh certification/refitting. Zero inverses, sign ties, singular systems and archive overflow have explicit rules. The bounds hold on every tape and include all unsuccessful attempts. The success proof uses only the conditional heavy-frequency coverage and the unchanged $P$ output cap of Lemma~\ref{lem:search}; its failure is at most $3\rho/4\le\rho$. Extras from light buckets face the same fresh original-target certificate as any inaccurate proposal. No additional all-searches-success event or independence among candidates is required. This completes Theorem~\ref{thm:main}.

\section{Scope, limitations and verification status}
The algorithm bounds both sample use and computation in the noiseless model of supplied examples. It does not improve Ye's logarithmic sample factor~\cite{Ye}, achieve the $O(k\log k)$ target, or provide a practical implementation. Its constants and polynomial exponents are deliberately loose. Noise robustness and further runtime improvements require separate analysis. The bound $k$ on target sparsity is an input promise. If the actual sparsity exceeds $k$, the finite resource caps still hold, but neither correctness nor detection of that violation is claimed. The final validation guarantee assumes a target of sparsity at most $k$.

The known $H$ can have degree $n$ and exponentially many Fourier coefficients. This is permitted because its evaluation cost is polynomial and it only proposes indices. The sparse residual used for estimation and certification is always $f-a$. The proof does not require a generic stability theorem for nonlinear corrections of sparse fits, or a short representation of every recovered index in the original samples.

After measurement, the algorithm is classical: its data consist of uniform labeled examples and successful Fourier samples with law $c_s^2$. The access model does not yield a general classical simulation of quantum computation or decide BPP versus BQP. Quantum-example preparation is not included in the theorem. Comparison with classical random-example bounds must keep this difference in supplied data explicit.

This manuscript has undergone internal AI-assisted mathematical, resource, attribution and artifact checking, with the review record separately identifying the reports for each manuscript version through Grok, Claude and AGY (configured with Gemini). It has not undergone independent peer review or formal certification. The accompanying review record identifies model provenance, findings, revisions and limits; document builds and exact arithmetic checks do not constitute a machine-checked proof.

\paragraph{Author and tool disclosure.}
Daniel Fredriksen is affiliated with Quantyra Inc. OpenAI Codex assisted the derivation, literature comparison, manuscript drafting and internal checking. Responsibility for the submitted claims remains with the author. The manuscript and accompanying archive are licensed under CC BY 4.0.

\begin{thebibliography}{9}
\bibitem{ACLPW}
S. Arunachalam, S. Chakraborty, T. Lee, M. Paraashar and R. de Wolf.
Two new results about quantum exact learning.
\emph{Quantum} 5, 587 (2021).
\href{https://doi.org/10.22331/q-2021-11-24-587}{doi:10.22331/q-2021-11-24-587}.
Corrected full version: \href{https://arxiv.org/abs/1810.00481v4}{arXiv:1810.00481v4}.
\bibitem{Ye}
Z. Ye. Nearly Optimal Quantum Sample Complexity of Learning Fourier-Sparse Boolean Functions.
\href{https://arxiv.org/abs/2610.17648v1}{arXiv:2610.17648v1} (2026).
\bibitem{Caro}
M. C. Caro, M. Hinsche, M. Ioannou, A. Nietner and R. Sweke.
Classical Verification of Quantum Learning.
\emph{ITCS 2024}, LIPIcs 287, 24:1--24:23.
\href{https://doi.org/10.4230/LIPIcs.ITCS.2024.24}{doi:10.4230/LIPIcs.ITCS.2024.24}.
Full version: \href{https://arxiv.org/abs/2306.04843v2}{arXiv:2306.04843v2} (2023).
\bibitem{GOSSW}
P. Gopalan, R. O'Donnell, R. A. Servedio, A. Shpilka and K. Wimmer.
Testing Fourier dimensionality and sparsity.
\emph{SIAM Journal on Computing} 40(4), 1075--1100 (2011).
\href{https://www.cs.columbia.edu/~rocco/Public/4-27-2011-gossw.pdf}{Author manuscript, April 27, 2011}.
\bibitem{KM}
E. Kushilevitz and Y. Mansour. Learning decision trees using the Fourier spectrum.
\emph{SIAM Journal on Computing} 22(6), 1331--1348 (1993).
\href{https://doi.org/10.1137/0222080}{doi:10.1137/0222080}.
\bibitem{Datta}
S. Datta, A. Ghosh, C. Kayal, M. Paraashar and M. Roy.
Spectral norm, economical sieve, and linear invariance testing of Boolean functions.
\emph{STACS 2026}, LIPIcs 364, 30:1--30:21 (2026).
\href{https://doi.org/10.4230/LIPIcs.STACS.2026.30}{doi:10.4230/LIPIcs.STACS.2026.30}.
\bibitem{Bareiss}
E. H. Bareiss. Sylvester's identity and multistep integer-preserving Gaussian elimination.
\emph{Mathematics of Computation} 22(103), 565--578 (1968).
\href{https://doi.org/10.1090/S0025-5718-1968-0226829-0}{doi:10.1090/S0025-5718-1968-0226829-0}.
\bibitem{Tropp}
J. A. Tropp. User-friendly tail bounds for sums of random matrices.
\emph{Foundations of Computational Mathematics} 12, 389--434 (2012).
\href{https://arxiv.org/abs/1004.4389v7}{arXiv:1004.4389v7} (2011).
\end{thebibliography}
\end{document}
